\section{Discussion}
\label{sec:discussion}

\paragraph{Optimality is context-dependent, and the context here is group size.}
The rule these agents use is optimal for the problem it was derived for: one
student, one teacher, a noisy channel whose reliability must be inferred online.
Nothing about it is careless. Yet the same rule, run in a large population with one
class-correlated perturbation, converts an informationally empty label into a
stable architecture of group distrust. The mechanism is specific: it acts through
$\Fmu$, and a rule with no affective sector has none to bias --- under a plain
Hebbian update the discrimination field does nothing at all, since the update does
not depend on $\hw$. We argue this rather than measure it: the ablation is cheap
and worth running, and we have not run it. What
exposes the population is precisely the machinery that makes each agent good at
inference: that it draws conclusions about \emph{sources} as well as about the
world. Giving interacting agents a well-founded model of each other's reliability
is not a neutral improvement.

\paragraph{Individually small, collectively discontinuous.} A single
discriminating agent is biased by $O(d)$; the population's state is not, moving
between phases over a narrow interval in $d$ once $\fd$ passes its threshold. An
evaluation protocol that inspects agents one at a time --- the standard protocol ---
can therefore certify every agent as approximately unbiased while the population
sits in a discriminatory phase. Auditing a multi-agent system requires
population-level order parameters, and the ones used here cost two elicited
matrices per society.

\paragraph{Two results that reframe standard questions.} The reversal at
$\alpha\approx1.7$ says that whether groups come to dislike each other because they
disagree, or the reverse, is not a choice between rival theories but two parameter
regimes of one system. And making $d$ negative, so that agents favour the
out-group, gives not a tidy inverted society but frustration:
$\BA<0$ with $\BI\approx0$, order parameters near zero not because nothing is
happening but because nothing can settle. Pair correlations alone would misreport
that state as neutral; only the balance measures identify it as a spin glass.

\paragraph{Limitations.} The agents are single-layer perceptrons over a fixed
shared embedding, all pairs may interact, and the society is closed and of fixed
size. The dynamics anneals rather than reaching a stationary state, so our phases
are phases at a finite interaction count. The class labels are exactly inert, the
clean version of a case that is never clean in practice. Most importantly, the
result is established on one substrate: agents whose learning rule we can write
down. Whether the same phases appear in a population whose agents are not
analysable --- language-model agents under instructions of varying tolerance, say
--- is the obvious next question, and the order parameters of
\S\ref{sec:order-params} were defined to make it answerable, but we do not answer
it here. Extending any of this to human societies is a further step again, which
the model does not license.
